%% kapa: exam=1 round=23 year=2012 period=1 form=A subject=economics subject_ko=경제원론 seq=1 number=41
%% source: Q-Net 감정평가사 기출문제 (cst00309) artlSeq=5052578 file=raw/1008/5052578/2012년 제23회 감정평가사 제1차시험 기출문제(1교시 A형).pdf
%% conversion: kordoc 4.24.51 -> tools/gen1.py (auto); review=proofread-sonnet; flags=none
\begin{question}{41}
완전경쟁시장에서 시장 수요함수가 $Q=1000-P$이고 기업들의 장기 평균비용은 생산량이 10일 때 100원으로 최소화된다. 이 때 장기균형에 관한 설명으로 옳지 않은 것은? (단, $Q$는 수요량, $P$는 가격이다.)
\begin{choices}
  \item 개별 기업의 이윤은 0원이다.
  \item 개별 기업의 생산량은 10이다.
  \item 균형가격은 100원이다.
  \item 시장에는 100개의 기업이 존재하게 된다.
  \item 소비자들은 가격순응자로서 효용을 극대화한다.
\end{choices}
\end{question}
\answer{4}
